\scanpage{28}
% archive sheet number: 27
dans $\unsure{\mathscr E|\mathcal E}$, qu'on peut choisir canoniquement en prenant le quart de grand cercle qui joint $P$ au milieu 2 du \deleted{segment} l'arc $(1,\infty)$ de grand cercle \unsure{opposée|opposé} à 0; ce qui définit un isom du rev universel de $\unsure{U|X}$ en $P$, \unsure{avec} celui de $U$ en $\unsure{\overline P|2}$ sur lequel le groupe précédent opère de façon canonique, d'où aussi par transport de structure sur $\widetilde U(\deleted{\unsure{\overline P}}P)$, donc il s'envoie dans $\mathscr E$.

On trouve que
\[
\dot\sigma_0\mapsto\sigma_0,\quad\dot\tau\to\tau_0
\]
est le remontage correspondant. De $=$ pour les pts fixes $-1$ de $\dot\sigma_1$, ou $\frac12$ de $\dot\sigma_\infty$, on trouve des scindages partiels sur les ss-groupes des \deleted{\unsure{groupes}} \unsure{stabs} $\Gamma_{-1}\subset\Gamma$ et $\Gamma_2\subset\Gamma$. Il est commode de poser
\[
\deleted{\unsure{2}}\quad 2=0',\qquad -1=1',\quad 1/2=\infty',
\]
d'où des groupes d'isotropie $\Gamma_{0'}=\{\dot\sigma_0,\tau\}$, $\Gamma_{1'}=\{\dot\sigma_1,\tau\}$, et $\Gamma_{\infty'}=\{\dot\sigma_\infty,\tau\}$, qui se remontent resp en $\sigma_0$ et $\tau_0$, en $\sigma_1$ et $\tau_1$, en $\sigma_\infty$ et $\tau_\infty$.

Les composantes connexes de $U^\tau$ sont les trois arcs de grand cercle entre $0,1,\infty$, qui contiennent chacun un \deleted{pt} et un seul des pts $0',1',\infty'$, et les classes de conjugaison
